% ભાગ: first-order-logic
% પ્રકરણ: beyond
% વિભાગ: intuitionistic-logic

\documentclass[../../../include/open-logic-section]{subfiles}

\begin{document}

\olfileid{fol}{byd}{il}

\olsection{સ્ફુરણાવાદી તર્કશાસ્ત્ર}

દ્વિતીય-ક્રમ અને ઉચ્ચ-ક્રમ તર્કશાસ્ત્રથી ઊલટું, સ્ફુરણાવાદી
પ્રથમ-ક્રમ તર્કશાસ્ત્ર પ્રશિષ્ટ સ્વરૂપનું એક મર્યાદન છે અને વધુ
``રચનાત્મક'' પ્રકારના તર્કવિચારનું નિદર્શન કરવા માટે બનાવવામાં આવ્યું
છે. નીચેના દાખલા તેની કેટલીક પાયાની પ્રેરણાઓ સ્પષ્ટ કરી શકે છે.

ધારો કે કોઈ એક દિવસ તમારી પાસે આવીને જાહેર કરે કે તેણે એવી પ્રાકૃતિક
સંખ્યા~$x$ નક્કી કરી છે કે જો~$x$ અવિભાજ્ય હોય તો રીમાન પરિકલ્પના સત્ય
છે, અને જો~$x$ સંયુક્ત હોય તો રીમાન પરિકલ્પના અસત્ય છે. બહુ સારા
સમાચાર!{} રીમાન પરિકલ્પના સત્ય છે કે નહિ એ ગણિતના મોટા વણઉકેલ્યા
પ્રશ્નોમાંનો એક છે, અને અહીં જાણે સમસ્યા ગણતરીના પ્રશ્નમાં, એટલે કોઈ
ચોક્કસ સંખ્યા અવિભાજ્ય છે કે નહિ તે નક્કી કરવામાં, ઘટી ગઈ છે.

$x$નું એ જાદુઈ મૂલ્ય શું છે? તેઓ તેનું વર્ણન આ રીતે કરે છે: રીમાન
પરિકલ્પના સત્ય હોય તો~$x$ બરાબર~$7$ અને નહિ તો~$9$ થાય એવી પ્રાકૃતિક
સંખ્યા છે.

ગુસ્સામાં તમે તમારા પૈસા પાછા માગો છો. પ્રશિષ્ટ દૃષ્ટિએ ઉપરનું વર્ણન
ખરેખર~$x$નું એકમાત્ર મૂલ્ય નક્કી કરે છે; પરંતુ તમને ખરેખર તો
\emph{સ્પષ્ટ રીતે} આપેલું~$x$નું મૂલ્ય જોઈએ છે.

બીજો, કદાચ થોડો ઓછો કૃત્રિમ, દાખલો જોઈએ. આપણને ખબર છે કે કોઈ અસંમેય
સંખ્યાને સંમેય ઘાતે ચઢાવવાથી સંમેય પરિણામ મળી શકે છે. દાખલા તરીકે,
$\sqrt{2}^2 = 2$. અસંમેય સંખ્યાને \emph{અસંમેય} ઘાતે ચઢાવીને સંમેય
પરિણામ મેળવી શકાય છે કે નહિ તે ઓછું સ્પષ્ટ છે. નીચેનો પ્રમેય તેનો
હકારાત્મક જવાબ આપે છે:

\begin{thm}
એવી અસંમેય સંખ્યાઓ $a$ અને~$b$ છે કે $a^b$ સંમેય છે.
\end{thm}

\begin{proof}
$\sqrt{2}^{\sqrt{2}}$ને વિચારો. આ સંખ્યા સંમેય હોય, તો આપણું કામ થઈ
ગયું: $a = b = \sqrt{2}$ લઈ શકીએ. નહિ તો તે અસંમેય છે. પછી
\[
(\sqrt{2}^{\sqrt{2}})^{\sqrt{2}} = \sqrt{2}^{\sqrt{2} \cdot
  \sqrt{2}} = \sqrt{2}^2 = 2,
\]
અને આ નિશ્ચિતપણે સંમેય છે. તેથી આ કિસ્સામાં $a$ને
$\sqrt{2}^{\sqrt{2}}$ અને $b$ને~$\sqrt 2$ લો.
\end{proof}

શું આ પ્રામાણ્ય સાબિતી છે? મોટા ભાગના ગણિતજ્ઞોને તે પ્રામાણ્ય લાગે છે.
પણ અહીં ફરી કંઈક થોડું અસંતોષકારક છે: કોઈ ચોક્કસ ગુણધર્મ ધરાવતા
વાસ્તવિક સંખ્યાઓના જોડાનું અસ્તિત્વ આપણે સાબિત કર્યું છે, પરંતુ તે
સંખ્યાઓનો \emph{કયો} જોડો છે તે કહી શકતા નથી. એ જ પરિણામ એવી રીતે પણ
સાબિત કરી શકાય છે કે જોડો $a$,~$b$ સાબિતીમાં જ \emph{આપેલો} હોય:
$a = \sqrt{3}$ અને $b = \log_3 4$ લો. પછી
\[
a^b = \sqrt{3}^{\log_3 4} = 3^{1/2 \cdot \log_3 4} = (3^{\log_3
  4})^{1/2} = 4^{1/2}= 2,
\]
કારણ કે $3^{\log_3 x} = x$.

સ્ફુરણાવાદી તર્કશાસ્ત્રમાં પહેલા પુરાવા જેવી ચાલ માન્ય ન હોય એવા
તર્કવિચારનું નિદર્શન કરવાનો આશય છે. $!A(x)$ સંતોષતા કોઈ~$x$નું
અસ્તિત્વ સાબિત કરવાનો અર્થ, બીજા પુરાવાની જેમ, ચોક્કસ~$x$ અને તે~$!A$
સંતોષે છે તેની સાબિતી આપવાનો છે. $!A$ અથવા~$!B$ સત્ય છે તે સાબિત કરવા
માટે બંનેમાંથી એકને સાબિત કરી શકવું જરૂરી છે.

વિધિવત્ રીતે કહીએ તો, પ્રથમ-ક્રમ તર્કશાસ્ત્રના !!a{derivation} તંત્રને
ચોક્કસ રીતે મર્યાદિત કરવાથી સ્ફુરણાવાદી પ્રથમ-ક્રમ તર્કશાસ્ત્ર મળે છે.
એ જ રીતે દ્વિતીય-ક્રમ અને ઉચ્ચ-ક્રમ તર્કશાસ્ત્રનાં પણ સ્ફુરણાવાદી
સ્વરૂપો છે. ગાણિતિક દૃષ્ટિએ આ માત્ર વિધિવત્ નિષ્પત્તિ-તંત્રો છે, પરંતુ
અગાઉ નોંધ્યું તેમ, તેમનો આશય એક પ્રકારના ગાણિતિક તર્કવિચારનું નિદર્શન
કરવાનો છે. તેને ગણિત વિશેના કોઈ ચોક્કસ તત્ત્વચિંતનાત્મક દૃષ્ટિકોણથી
ન્યાયસંગત તર્કવિચાર (જેમ કે બ્રાઉરનો સ્ફુરણાવાદ), અથવા વધુ ``મૂર્ત''
અને સંતોષકારક ગાણિતિક તર્કવિચાર (બિશપના રચનાવાદની દિશામાં) માની શકાય;
અને આ વિધિવત્ વર્ણન અનૌપચારિક પ્રેરણાને પકડે છે કે નહિ તે અંગે દલીલ
કરી શકાય. પરંતુ આપણું તત્ત્વચિંતનાત્મક વલણ જે હોય તે, સ્ફુરણાવાદી
તર્કશાસ્ત્રનો વિધિવત્ રીતે રજૂ થયેલા તર્કશાસ્ત્ર તરીકે અભ્યાસ કરી
શકીએ છીએ; અને કારણ ગમે તે હોય, ઘણા ગાણિતિક તર્કશાસ્ત્રીઓને તેનો
અભ્યાસ રસપ્રદ લાગે છે.

સ્ફુરણાવાદી સંયોજકોનું એક અનૌપચારિક રચનાત્મક અર્થઘટન છે, જેને સામાન્ય
રીતે BHK અર્થઘટન (બ્રાઉર, હાઇટિંગ અને કોલ્મોગોરોવ પરથી નામ આપેલું)
કહે છે. તે આ પ્રમાણે છે: $!A \land !B$ની સાબિતી એટલે~$!A$ની સાબિતી
સાથે જોડેલી~$!B$ની સાબિતી; $!A \lor !B$ની સાબિતી એટલે~$!A$ની અથવા
$!B$ની સાબિતી, અને કયો કિસ્સો છે તેની સ્પષ્ટ માહિતી; $!A \lif !B$ની
સાબિતી એટલે~$!A$ની સાબિતીને~$!B$ની સાબિતીમાં ફેરવતી પ્રક્રિયા;
$\lforall[x][!A(x)]$ની સાબિતી એટલે~$x$ના કોઈ પણ મૂલ્ય માટે~$!A(x)$ની
સાબિતી આપતી પ્રક્રિયા; અને $\lexists[x][!A(x)]$ની સાબિતી એટલે~$x$નું
કોઈ મૂલ્ય અને તે મૂલ્ય~$!A$ સંતોષે છે તેની સાબિતી. આ અર્થઘટનને
``કરી--હાવર્ડ એકરૂપતા'' અથવા ``!!{formula}s-રૂપે-પ્રકારો પરિરૂપ'' નામે
ઓળખાતા ગણનાત્મક પદોમાં વર્ણવી શકાય છે: !!a{formula}ને કોઈ ચોક્કસ
માહિતી-પ્રકાર નિર્દેશતું માનો અને સાબિતીઓને આ માહિતી-પ્રકારોની એવી
ગણનાત્મક વસ્તુઓ માનો જે આપણને અનુરૂપ !!{formula} સત્ય છે તે જોવા દે છે.

સ્ફુરણાવાદી તર્કશાસ્ત્રને ઘણી વાર પ્રશિષ્ટ તર્કશાસ્ત્રમાંથી વર્જિત
મધ્યનો નિયમ ``બાદ'' કરવાથી મળેલું માનવામાં આવે છે. નીચેનો પ્રમેય આ
વાતને વધુ ચોક્કસ બનાવે છે.
\begin{thm}
સ્ફુરણાવાદી રીતે નીચેના સ્વયંસિદ્ધિ-પ્રરૂપો સમકક્ષ છે:
\begin{enumerate}
\item $(\lnot !A \lif \lfalse) \lif !A$.
\item $!A \lor \lnot !A$
\item $\lnot \lnot !A \lif !A$
\end{enumerate}
\end{thm}
કોઈ એક પ્રરૂપનાં દાખલાઓ બીજા બેમાંથી કોઈ એકમાંથી મેળવવા એ સ્ફુરણાવાદી
તર્કશાસ્ત્રનો સારો સ્વાધ્યાય છે.

બ્રાઉરના સ્ફુરણાવાદનાં વિધિવતીકરણ તરીકે રજૂ થયેલાં સ્ફુરણાવાદી
વિધાનાત્મક તર્કશાસ્ત્રનાં પ્રથમ નિષ્પત્તિ-તંત્રો કોલ્મોગોરોવ, ગ્લિવેન્કો
અને હાઇટિંગે સ્વતંત્ર રીતે આપ્યાં હતાં. સ્ફુરણાવાદી પ્રથમ-ક્રમ
તર્કશાસ્ત્રનું (અને સ્ફુરણાવાદી ગણિતના કેટલાક ભાગોનું) પહેલું
વિધિવતીકરણ હાઇટિંગે આપ્યું હતું. પ્રશિષ્ટ રીતે પ્રામાણ્ય ધરાવતા ઘણાં
પ્રરૂપો સ્ફુરણાવાદી રીતે પ્રામાણ્ય ધરાવતા નથી, છતાં ઘણાં ધરાવે છે.

\emph{દ્વિનિષેધ અનુવાદ} પ્રશિષ્ટ અને સ્ફુરણાવાદી તર્કશાસ્ત્ર વચ્ચેનો
મહત્વનો સંબંધ વર્ણવે છે. તે નીચે મુજબ અનુમાનાત્મક રીતે વ્યાખ્યાયિત થાય
છે ($!A^N$ને પ્રશિષ્ટ !!{formula}~$!A$નો ``સ્ફુરણાવાદી'' અનુવાદ માનો):
\begin{align*}
!A^N & \ident \lnot\lnot !A \quad \text{આણ્વિક !!{formula}s $!A$ માટે} \\
(!A \land !B)^N & \ident (!A^N \land !B^N) \\
(!A \lor !B)^N & \ident  \lnot\lnot (!A^N \lor !B^N) \\
(!A \lif !B)^N & \ident (!A^N \lif !B^N) \\
(\lforall[x][!A])^N & \ident \lforall[x][!A^N] \\
(\lexists[x][!A])^N & \ident \lnot\lnot\lexists[x][!A^N]
\end{align*}
વિધાનાત્મક તર્કશાસ્ત્ર માટે કોલ્મોગોરોવ અને ગ્લિવેન્કો પાસે આ અનુવાદનાં
સ્વરૂપો હતાં; વિધેય તર્કશાસ્ત્ર માટે તે ગડેલ અને ગેન્ટ્સને સ્વતંત્ર
રીતે આપ્યો હતો. આપણી પાસે નીચેનું છે:

\begin{thm}
\begin{enumerate}
\item $!A \liff !A^N$ પ્રશિષ્ટ રીતે સાબિત કરી શકાય છે.
\item જો~$!A$ પ્રશિષ્ટ રીતે સાબિત કરી શકાય, તો $!A^N$ સ્ફુરણાવાદી રીતે
  સાબિત કરી શકાય છે.
\end{enumerate}
\end{thm}

હવે નીચેના સંવાદની કલ્પના કરી શકીએ. પ્રશિષ્ટ ગણિતજ્ઞ: ``મેં~$!A$
સાબિત કર્યું!'' સ્ફુરણાવાદી ગણિતજ્ઞ: ``તમારી સાબિતી પ્રામાણ્ય નથી.
તમે ખરેખર તો~$!A^N$ સાબિત કર્યું છે.'' પ્રશિષ્ટ ગણિતજ્ઞ: ``મને એ પણ
ચાલે!'' પ્રશિષ્ટ ગણિતજ્ઞની દૃષ્ટિએ સ્ફુરણાવાદી ગણિતજ્ઞ ઝીણી બાબતનો
અતિશય આગ્રહ રાખે છે, કારણ કે બંને સમકક્ષ છે. પરંતુ સ્ફુરણાવાદી ગણિતજ્ઞ
માને છે કે બંનેમાં તફાવત છે.

નોંધો કે ઉપરનો અનુવાદ માત્ર શુદ્ધ તર્કશાસ્ત્રને લગતો છે; પ્રશિષ્ટ અને
સ્ફુરણાવાદી ગણિત માટે યોગ્ય \emph{અતાર્કિક} સ્વયંસિદ્ધિઓ કઈ છે અથવા
તેમની વચ્ચે શું સંબંધ છે તે પ્રશ્નનો તે જવાબ આપતો નથી. પરંતુ ઉપરના
પ્રમેયનો નીચેનો નાનો વિસ્તાર કેટલીક ઉપયોગી માહિતી આપે છે:

\begin{thm}
જો~$\Gamma$,~$!A$ને પ્રશિષ્ટ રીતે સાબિત કરે, તો $\Gamma^N$,~$!A^N$ને
સ્ફુરણાવાદી રીતે સાબિત કરે છે.
\end{thm}

બીજા શબ્દોમાં, કેટલીક ધારણાઓમાંથી~$!A$ પ્રશિષ્ટ રીતે સાબિત કરી શકાય,
તો તેમના દ્વિનિષેધ અનુવાદોમાંથી~$!A^N$ સ્ફુરણાવાદી રીતે સાબિત કરી શકાય
છે.

કોઈ વાક્ય અથવા વિધાનાત્મક !!{formula} સ્ફુરણાવાદી રીતે પ્રામાણ્ય ધરાવે
છે તે બતાવવા માટે માત્ર એક સાબિતી આપવી પડે છે. પરંતુ તે પ્રામાણ્ય
ધરાવતું નથી એ કેવી રીતે બતાવશો? તે માટે આપણને યથાર્થ અને શક્ય હોય તો
પૂર્ણ અર્થવિજ્ઞાન જોઈએ. ક્રિપ્કેએ આપેલું અર્થવિજ્ઞાન આ કામ માટે સરસ
રીતે બંધબેસે છે.

પ્રશિષ્ટ તર્કશાસ્ત્ર માટે કરેલી જ પ્રક્રિયા અહીં કરી શકીએ: અર્થવિજ્ઞાન
વ્યાખ્યાયિત કરી તેના માટે યથાર્થતા અને પૂર્ણતા સાબિત કરી શકીએ. જોકે
નીચેનો તફાવત નોંધવા જેવો છે. પ્રશિષ્ટ તર્કશાસ્ત્રના કિસ્સામાં
અર્થવિજ્ઞાન એક અર્થમાં ``સ્વાભાવિક'' હતું અને સંયોજકોના અર્થમાં
અંતર્નિહિત હતું. ક્રિપ્કે અર્થવિજ્ઞાન માટે થોડી સહજ પ્રેરણા આપી શકાય,
પણ તેમાં એવી અનિવાર્યતાનો અનુભવ થતો નથી. ઉપરાંત, પ્રશિષ્ટ
!!{structure}નો ખ્યાલ કુદરતી ગાણિતિક ખ્યાલ છે; તેથી !!a{structure}ના
ખ્યાલને પ્રશિષ્ટ પ્રથમ-ક્રમ તર્કશાસ્ત્રના અભ્યાસનું સાધન માની શકાય,
અથવા પ્રશિષ્ટ પ્રથમ-ક્રમ તર્કશાસ્ત્રને ગાણિતિક !!{structure}sના
અભ્યાસનું સાધન માની શકાય. તેનાથી ઊલટું, ક્રિપ્કે !!{structure}sને માત્ર
તાર્કિક રચના તરીકે જ જોઈ શકાય; તેમનું અલગથી કોઈ ગાણિતિક મહત્ત્વ હોય
એવું લાગતું નથી.

વિધાનાત્મક ભાષા માટેની ક્રિપ્કે !!{structure}~$\mModel M = \tuple{W,
R, V}$માં ગણ~$W$, લઘુતમ !!{element} ધરાવતો~$W$ પરનો આંશિક ક્રમ~$R$,
અને વિધાનાત્મક ચલોનું~$W$ના !!{element}sને એક ``એકદિશવર્ધી'' નિયુક્તિ
હોય છે. સહજ વિચાર એવો છે કે~$W$ના !!{element}s ``જગતો'' અથવા
``જ્ઞાનની અવસ્થાઓ''નું પ્રતિનિધિત્વ કરે છે; ઘટક $v \geq u$,
~$u$ની ``શક્ય ભાવિ અવસ્થા''નું પ્રતિનિધિત્વ કરે છે; અને~$u$ને નિયુક્ત
વિધાનાત્મક ચલો એ અવસ્થા~$u$માં સત્ય હોવાનું જાણીતું હોય તે વિધાનો છે.
પછી બાધન-સંબંધ $\mSat{M}{!A}[w]$ આ સંબંધને ભાષાના મનસ્વી
!!{formula}s સુધી વિસ્તારે છે; $\mSat{M}{!A}[w]$ને ``અવસ્થા~$w$માં
$!A$ સત્ય છે'' એમ વાંચો. સંબંધ નીચે મુજબ અનુમાનાત્મક રીતે વ્યાખ્યાયિત
થાય છે:
\begin{enumerate}
\item $\mSat{M}{\Obj p_i}[w]$ ત્યારે અને માત્ર ત્યારે જ્યારે $\Obj p_i$,
  ~ $w$ને નિયુક્ત વિધાનાત્મક ચલોમાંનો એક હોય.
\item $\mSat/{M}{\lfalse}[w]$.
\item $\mSat{M}{(!A \land !B)}[w]$ ત્યારે અને માત્ર ત્યારે જ્યારે
  $\mSat{M}{!A}[w]$ અને $\mSat{M}{!B}[w]$.
\item $\mSat{M}{(!A \lor !B)}[w]$ ત્યારે અને માત્ર ત્યારે જ્યારે
  $\mSat{M}{!A}[w]$ અથવા $\mSat{M}{!B}[w]$.
\item $\mSat{M}{(!A \lif !B)}[w]$ ત્યારે અને માત્ર ત્યારે જ્યારે, દરેક
  $w' \geq w$ માટે $\mSat{M}{!A}[w']$ હોય ત્યારે $\mSat{M}{!B}[w']$.
\end{enumerate}
$\lnot (p \land q) \lif (\lnot p \lor \lnot q)$ સ્ફુરણાવાદી રીતે
પ્રામાણ્ય ધરાવતું નથી તે બતાવતી પ્રતિદૃષ્ટાંતરૂપ ક્રિપ્કે !!{structure}
રચવાનો પ્રયત્ન કરવો સારો સ્વાધ્યાય છે.

\end{document}
